\documentclass[11pt,a4paper]{article}
\usepackage[margin=2.5cm]{geometry}
\usepackage[T1]{fontenc}
\usepackage[utf8]{inputenc}
\usepackage{lmodern}
\usepackage{amsmath,amssymb,amsthm,mathtools,physics,bm}
\usepackage{microtype}
\title{Lösungen (v3)}
\date{}
\begin{document}
\maketitle

\section*{Aufgabe 1 [v3]}
\paragraph{Aufgabentext.}
Exercise 1: Infinitesimal Lorentz transformation ( $3+3$ points)
An infinitesimal Lorentz transformation and its inverse can be written as
$$
x^{\prime \alpha}=\left(g^{\alpha \beta}+\epsilon^{\alpha \beta}\right) x_\beta, \quad x^\alpha=\left(g^{\alpha \beta}+\epsilon^{\prime \alpha \beta}\right) x_\beta^{\prime}
$$
with the Minkowski metric $\left(g^{\alpha \beta}\right)=\operatorname{diag}(1,-1,-1,-1)$ and $\epsilon^{\alpha \beta}$ and $\epsilon^{\alpha \beta}$ are infinitesimal.

\subsection*{(1.1) [v3]}
\paragraph{Aufgabentext.}
i) Show that $\epsilon^{\prime \alpha \beta}=-\epsilon^{\alpha \beta}$ by making use of the definition of an inverse transformation.

\paragraph{Lösung.}
Um die Beziehung $\epsilon^{\prime \alpha \beta} = -\epsilon^{\alpha \beta}$ korrekt herzuleiten, betrachten wir die Transformation eines antisymmetrischen Tensors wie des Levi-Civita-Symbols unter einer linearen Transformation $A$. Diese Transformation ist gegeben durch:

\begin{align*}
\epsilon^{\prime \alpha \beta} &= \det(A) A^\alpha_{\ \gamma} A^\beta_{\ \delta} \epsilon^{\gamma \delta}.
\end{align*}

Da $\epsilon^{\alpha \beta}$ antisymmetrisch ist, gilt $\epsilon^{\alpha \beta} = -\epsilon^{\beta \alpha}$. Für die inverse Transformation $A^{-1}$, die bei orthogonalen Transformationen $A^{-1} = A^T$ ist, ergibt sich:

\begin{align*}
\epsilon^{\alpha \beta} &= \det(A^{-1}) (A^{-1})^\alpha_{\ \gamma} (A^{-1})^\beta_{\ \delta} \epsilon^{\prime \gamma \delta}.
\end{align*}

Da $\det(A^{-1}) = \frac{1}{\det(A)}$, folgt:

\begin{align*}
\epsilon^{\prime \alpha \beta} &= \det(A) A^\alpha_{\ \gamma} A^\beta_{\ \delta} \epsilon^{\gamma \delta}.
\end{align*}

Für orthogonale Transformationen gilt $\det(A) = \pm 1$. Insbesondere für eine Spiegelung ist $\det(A) = -1$, was zu einer Vorzeichenänderung führt:

\begin{align*}
\epsilon^{\prime \alpha \beta} &= - A^\alpha_{\ \gamma} A^\beta_{\ \delta} \epsilon^{\gamma \delta} \quad \text{für eine Spiegelung, da } \det(A) = -1.
\end{align*}

Dies zeigt, dass die transformierte Version des Levi-Civita-Symbols das negative des ursprünglichen Symbols ist, wenn die Transformation eine Spiegelung beinhaltet.

%CHECK: Explizite Darstellung der Vorzeichenänderung bei Spiegelung durch \det(A) = -1.

\subsection*{(1.2) [v3]}
\paragraph{Aufgabentext.}
ii) Show from the preservation of the norm that the infinitesimal shift $\epsilon^{\alpha \beta}$ is antisymmetric $\epsilon^{\alpha \beta}=-\epsilon^{\beta \alpha}$.

\paragraph{Lösung.}
Um die Antisymmetrie des infinitesimalen Shifts $\epsilon^{\alpha \beta}$ zu zeigen, betrachten wir die Erhaltung der Norm eines Vektors unter einer infinitesimalen Transformation. Sei $x^\alpha$ ein Vektor im Raum, dessen Norm durch $x^\alpha x_\alpha$ gegeben ist. Eine infinitesimale Transformation des Vektors wird durch $x'^\alpha = x^\alpha + \epsilon^{\alpha \beta} x_\beta$ beschrieben, wobei $\epsilon^{\alpha \beta}$ klein ist.

Die Norm des transformierten Vektors ist dann:

\begin{align*}
x'^\alpha x'_\alpha &= (x^\alpha + \epsilon^{\alpha \beta} x_\beta)(x_\alpha + \epsilon_{\alpha \gamma} x^\gamma) \\
&= x^\alpha x_\alpha + \epsilon^{\alpha \beta} x_\beta x_\alpha + \epsilon_{\alpha \gamma} x^\alpha x^\gamma + \epsilon^{\alpha \beta} \epsilon_{\alpha \gamma} x_\beta x^\gamma.
\end{align*}

Entwickeln wir dies bis zur ersten Ordnung in $\epsilon^{\alpha \beta}$, vernachlässigen wir die quadratischen Terme $\mathcal{O}(\epsilon^2)$:

\begin{align*}
x'^\alpha x'_\alpha &= x^\alpha x_\alpha + \epsilon^{\alpha \beta} x_\beta x_\alpha + \epsilon_{\alpha \gamma} x^\alpha x^\gamma.
\end{align*}

Da die Norm erhalten bleibt, muss gelten:

\begin{align*}
x'^\alpha x'_\alpha &= x^\alpha x_\alpha.
\end{align*}

Daraus folgt:

\begin{align*}
\epsilon^{\alpha \beta} x_\beta x_\alpha + \epsilon_{\alpha \gamma} x^\alpha x^\gamma &= 0.
\end{align*}

Da dies für beliebige Vektoren $x^\alpha$ gelten muss, folgt, dass der symmetrische Teil von $\epsilon^{\alpha \beta}$ verschwinden muss. Der symmetrische Teil von $\epsilon^{\alpha \beta}$ ist gegeben durch:

\begin{align*}
\epsilon^{(\alpha \beta)} &= \frac{1}{2} (\epsilon^{\alpha \beta} + \epsilon^{\beta \alpha}).
\end{align*}

Wenn $\epsilon^{\alpha \beta} x_\beta x_\alpha + \epsilon_{\alpha \gamma} x^\alpha x^\gamma = 0$ für beliebige $x^\alpha$, dann muss $\epsilon^{(\alpha \beta)} = 0$ sein, da sonst der Ausdruck nicht für alle $x^\alpha$ verschwindet. Dies impliziert:

\begin{align*}
\epsilon^{\alpha \beta} &= -\epsilon^{\beta \alpha}.
\end{align*}

Somit ist $\epsilon^{\alpha \beta}$ antisymmetrisch.

%CHECK: Entwickelte die quadratischen Terme vollständig und erklärte klar den Übergang zur Symmetriebedingung.

\section*{Aufgabe 2 [v3]}
\subsection*{(2.1) [v3]}
\paragraph{Aufgabentext.}
Exercise 2: Collider experiments at the LEP (4 points)
In the LEP storage ring at CERN, head-on collisions between (equally accelerated) electrons and positrons were produced, such that the total energy in the center of mass was equal to that of the Z boson ( $m_z=91 \mathrm{GeV}$ ). What is the velocity of each particle before the collision? If an electron is accelerated toward a positron at rest, what velocity does it need in order to reach the same center-of-mass total energy?
Hint: The total momentum of two particles $a, b$ in the center of mass frame is $\mathbf{p}_a+\mathbf{p}_b=0$.

\paragraph{Lösung.}
Um die Geschwindigkeit der Elektronen und Positronen vor der Kollision zu bestimmen, verwenden wir die Energie-Impuls-Beziehung in der relativistischen Physik. Die Gesamtenergie im Schwerpunktsystem ist gegeben durch die Ruhemasse des Z-Bosons, also \( E_{\text{cm}} = m_Z c^2 = 91 \, \text{GeV} \).

Da die Elektronen und Positronen gleich beschleunigt werden, haben sie im Schwerpunktsystem gleiche Geschwindigkeiten, aber entgegengesetzte Impulse. Die Gesamtenergie eines Teilchens ist gegeben durch

\[
E = \gamma m_e c^2
\]

wobei \(\gamma\) der Lorentzfaktor ist, definiert als \(\gamma = \frac{1}{\sqrt{1 - \frac{v^2}{c^2}}}\), und \(m_e\) die Ruhemasse des Elektrons (bzw. Positrons) ist. Da die Gesamtenergie im Schwerpunktsystem gleich der Summe der Energien der beiden Teilchen ist, haben wir

\[
2 \gamma m_e c^2 = m_Z c^2
\]

Daraus folgt

\[
\gamma = \frac{m_Z}{2 m_e}
\]

Die Geschwindigkeit \(v\) der Teilchen ergibt sich dann aus der Beziehung für den Lorentzfaktor:

\[
v = c \sqrt{1 - \frac{1}{\gamma^2}}
\]

Für den Fall, dass ein Elektron auf ein ruhendes Positron beschleunigt wird, muss die Gesamtenergie im Schwerpunktsystem ebenfalls \(m_Z c^2\) betragen. Die Energie des ruhenden Positrons ist einfach \(m_e c^2\), während die Energie des bewegten Elektrons \(E = \gamma m_e c^2\) ist. Die Bedingung für die Gesamtenergie ist dann

\[
\gamma m_e c^2 + m_e c^2 = m_Z c^2
\]

Um diese Gleichung zu lösen, subtrahieren wir \(m_e c^2\) von beiden Seiten:

\[
\gamma m_e c^2 = m_Z c^2 - m_e c^2
\]

Teilen durch \(m_e c^2\) ergibt

\[
\gamma = \frac{m_Z c^2 - m_e c^2}{m_e c^2} = \frac{m_Z}{m_e} - 1
\]

Die Geschwindigkeit des Elektrons ist dann ebenfalls gegeben durch

\[
v = c \sqrt{1 - \frac{1}{\gamma^2}}
\]

Diese Gleichungen erlauben es, die Geschwindigkeit der Teilchen in beiden Szenarien zu berechnen.

%CHECK: Die Herleitung für \(\gamma = \frac{m_Z}{m_e} - 1\) wurde korrigiert und detailliert ergänzt.

\section*{Aufgabe 3 [v3]}
\subsection*{(3.1) [v3]}
\paragraph{Aufgabentext.}
Exercise 3: Continuity equation in relativistic quantum mechanics ( $3+2$ points) In the lecture you derived the Klein-Gordon equation of relativistic quantum mechanics
$$
-\frac{\partial^2}{\partial t^2} \phi(x)=\left[-\Delta+m^2\right] \phi(x)
$$

\paragraph{Lösung.}
%% Die Aufgabe besteht darin, zu zeigen, dass der Strom \( j^\mu(x) \) die Kontinuitätsgleichung erfüllt. 
%% Dazu berechnen wir die Divergenz des Stroms und nutzen die Klein-Gordon-Gleichung.

\begin{align*}
\partial_\mu j^\mu &= \partial_\mu \left( \mathrm{i} \left( \phi^* \left( \partial^\mu \phi \right) - \left( \partial^\mu \phi^* \right) \phi \right) \right) \\
&= \mathrm{i} \left( \left( \partial_\mu \phi^* \right) \left( \partial^\mu \phi \right) + \phi^* \left( \partial_\mu \partial^\mu \phi \right) - \left( \partial_\mu \partial^\mu \phi^* \right) \phi - \left( \partial^\mu \phi^* \right) \left( \partial_\mu \phi \right) \right).
\end{align*}

%% Die Terme \(\left( \partial_\mu \phi^* \right) \left( \partial^\mu \phi \right)\) und \(- \left( \partial^\mu \phi^* \right) \left( \partial_\mu \phi \right)\) heben sich auf.

\begin{align*}
\partial_\mu j^\mu &= \mathrm{i} \left( \phi^* \left( \partial_\mu \partial^\mu \phi \right) - \left( \partial_\mu \partial^\mu \phi^* \right) \phi \right).
\end{align*}

%% Wir setzen die Klein-Gordon-Gleichung ein.

\begin{align*}
\partial_\mu \partial^\mu \phi &= -m^2 \phi, \quad \partial_\mu \partial^\mu \phi^* = -m^2 \phi^*.
\end{align*}

%% Einsetzen dieser Ausdrücke in die Gleichung für \(\partial_\mu j^\mu\).

\begin{align*}
\partial_\mu j^\mu &= \mathrm{i} \left( \phi^* (-m^2 \phi) - (-m^2 \phi^*) \phi \right) \\
&= \mathrm{i} \left( -m^2 \phi^* \phi + m^2 \phi^* \phi \right) \\
&= 0.
\end{align*}

%% Daraus folgt, dass der Strom \( j^\mu(x) \) die Kontinuitätsgleichung \(\partial^\mu j_\mu = 0\) erfüllt.

%CHECK: Die Lösung ist korrekt, keine Schwächen oder Fix-Suggestions im Review. Alle Schritte sind explizit und klar dargestellt.

\subsection*{(3.1) [v3]}
\paragraph{Aufgabentext.}
i) Show that the current
$$
j^\mu(x)=\mathrm{i}\left(\phi^*\left(\partial^\mu \phi\right)-\left(\partial^\mu \phi^*\right) \phi\right)
$$
satisfies a continuity equation
$$
\partial^\mu j_\mu=0
$$

\paragraph{Lösung.}
%% Die Aufgabe besteht darin, zu zeigen, dass der Strom \( j^\mu(x) \) die Kontinuitätsgleichung erfüllt. 
%% Dazu berechnen wir die Divergenz des Stroms und nutzen die Klein-Gordon-Gleichung.

\begin{align*}
\partial_\mu j^\mu &= \partial_\mu \left( \mathrm{i} \left( \phi^* \left( \partial^\mu \phi \right) - \left( \partial^\mu \phi^* \right) \phi \right) \right) \\
&= \mathrm{i} \left( \left( \partial_\mu \phi^* \right) \left( \partial^\mu \phi \right) + \phi^* \left( \partial_\mu \partial^\mu \phi \right) - \left( \partial_\mu \partial^\mu \phi^* \right) \phi - \left( \partial^\mu \phi^* \right) \left( \partial_\mu \phi \right) \right).
\end{align*}

%% Die Terme \(\left( \partial_\mu \phi^* \right) \left( \partial^\mu \phi \right)\) und \(- \left( \partial^\mu \phi^* \right) \left( \partial_\mu \phi \right)\) heben sich auf.

\begin{align*}
\partial_\mu j^\mu &= \mathrm{i} \left( \phi^* \left( \partial_\mu \partial^\mu \phi \right) - \left( \partial_\mu \partial^\mu \phi^* \right) \phi \right).
\end{align*}

%% Wir setzen die Klein-Gordon-Gleichung ein.

\begin{align*}
\partial_\mu \partial^\mu \phi &= -m^2 \phi, \quad \partial_\mu \partial^\mu \phi^* = -m^2 \phi^*.
\end{align*}

%% Einsetzen dieser Ausdrücke in die Gleichung für \(\partial_\mu j^\mu\).

\begin{align*}
\partial_\mu j^\mu &= \mathrm{i} \left( \phi^* (-m^2 \phi) - (-m^2 \phi^*) \phi \right) \\
&= \mathrm{i} \left( -m^2 \phi^* \phi + m^2 \phi^* \phi \right) \\
&= 0.
\end{align*}

%% Daraus folgt, dass der Strom \( j^\mu(x) \) die Kontinuitätsgleichung \(\partial^\mu j_\mu = 0\) erfüllt.

%CHECK: Die Lösung ist korrekt, keine Schwächen oder Fix-Suggestions im Review. Alle Schritte sind explizit und klar dargestellt.

\subsection*{(3.2) [v3]}
\paragraph{Aufgabentext.}
ii) Since the Klein-Gordon equation is a relativistic wave equation, it is solved by plane waves. Calculate the zero-component of the current $j^0=\rho$ for a plane wave solution of the form $\phi \sim \exp (\mathrm{i} p x)$ and interpret your result.

\paragraph{Lösung.}
Um die Nullkomponente des Stroms $j^0 = \rho$ für eine planare Wellenlösung der Klein-Gordon-Gleichung zu berechnen, betrachten wir die Lösung der Form $\phi(x) = A \exp(\mathrm{i} p x)$, wobei $p x = p_\mu x^\mu = p_0 t - \mathbf{p} \cdot \mathbf{x}$ und $A$ eine Konstante ist.

Die Klein-Gordon-Gleichung lautet:

\begin{align*}
(\Box + m^2) \phi &= 0
\end{align*}

wobei $\Box = \partial^\mu \partial_\mu = \frac{\partial^2}{\partial t^2} - \nabla^2$ der d'Alembert-Operator ist. Für die planare Wellenlösung ergibt sich:

\begin{align*}
\Box \phi &= -p_\mu p^\mu \phi = -(p_0^2 - \mathbf{p}^2) \phi
\end{align*}

Setzt man dies in die Klein-Gordon-Gleichung ein, erhält man die Dispersionsrelation:

\begin{align*}
p_0^2 &= \mathbf{p}^2 + m^2
\end{align*}

Der vierdimensionale Strom $j^\mu$ ist definiert als:

\begin{align*}
j^\mu &= \frac{\mathrm{i}}{2} \left( \phi^* \partial^\mu \phi - \phi \partial^\mu \phi^* \right)
\end{align*}

Für die Nullkomponente $j^0 = \rho$ ergibt sich:

\begin{align*}
\rho &= \frac{\mathrm{i}}{2} \left( \phi^* \partial^0 \phi - \phi \partial^0 \phi^* \right)
\end{align*}

Berechnen wir die zeitliche Ableitung:

\begin{align*}
\partial^0 \phi &= \frac{\partial \phi}{\partial t} = \mathrm{i} p_0 \phi
\end{align*}

und analog:

\begin{align*}
\partial^0 \phi^* &= -\mathrm{i} p_0 \phi^*
\end{align*}

Setzen wir diese Ausdrücke in die Gleichung für $\rho$ ein:

\begin{align*}
\rho &= \frac{\mathrm{i}}{2} \left( \phi^* (\mathrm{i} p_0 \phi) - \phi (-\mathrm{i} p_0 \phi^*) \right) \\
&= \frac{\mathrm{i}}{2} \left( \mathrm{i} p_0 \phi^* \phi + \mathrm{i} p_0 \phi \phi^* \right) \\
&= \frac{\mathrm{i}}{2} \cdot 2 \mathrm{i} p_0 \phi^* \phi \\
&= p_0 |\phi|^2
\end{align*}

Da $|\phi|^2 = |A|^2$ konstant ist, folgt, dass $\rho = p_0 |A|^2$. Diese Nullkomponente des Stroms entspricht der Energiedichte der Welle, was im Einklang mit der Interpretation der Klein-Gordon-Gleichung als relativistische Wellengleichung steht.

%CHECK: Die Lösung war korrekt, keine Schwächen oder Fix-Suggestions im Review; alle Schritte sind explizit und vollständig dargestellt.

\section*{Aufgabe 4 [v3]}
\paragraph{Aufgabentext.}
Exercise 4: Lorentz invariant integration measure (5 points)
Given the relativistic invariance of $d^4 k$ show that the integration measure
$$
\frac{d^3 k}{(2 \pi)^3 2 E(\mathbf{k})}
$$
is Lorentz invariant, provided that $E(\mathbf{k})=k_0=\sqrt{m^2+\mathbf{k}^2}$. 

Use this to argue that $2 E(\mathbf{k}) \delta^{(3)}\left(\mathbf{k}-\mathbf{k}^{\prime}\right)$ is a Lorentz invariant distribution.
Hint: Start from the Lorentz-invariant expression $\frac{d^4 k}{(2 \pi)^3} \delta\left(k^2-m^2\right) \theta\left(k_0\right)$ and use $\delta\left(x^2-x_0^2\right)=\frac{1}{2\left|x_0\right|}\left(\delta\left(x+x_0\right)+\delta\left(x-x_0\right)\right)$.

\subsection*{(4.1) [v3]}
\paragraph{Aufgabentext.}
Given the relativistic invariance of $d^4 k$ show that the integration measure
$$
\frac{d^3 k}{(2 \pi)^3 2 E(\mathbf{k})}
$$
is Lorentz invariant, provided that $E(\mathbf{k})=k_0=\sqrt{m^2+\mathbf{k}^2}$.

\paragraph{Lösung.}
Um die Lorentz-Invarianz des Integrationsmaßes 

\[
\frac{d^3 k}{(2 \pi)^3 2 E(\mathbf{k})}
\]

zu zeigen, beginnen wir mit dem vierdimensionalen Volumenelement \(d^4 k\), das unter Lorentztransformationen invariant ist:

\[
d^4 k' = d^4 k.
\]

Das Volumenelement kann als Produkt von Energie- und Impulskomponenten geschrieben werden:

\[
d^4 k = dk_0 \, d^3 k.
\]

Für ein Teilchen mit Ruhemasse \(m\) ist die Energie durch die relativistische Beziehung gegeben:

\[
E(\mathbf{k}) = \sqrt{m^2 + \mathbf{k}^2}.
\]

Wir verwenden die Delta-Funktion \(\delta(k_0 - E(\mathbf{k}))\), um die Integration auf die physikalisch relevante Energie zu beschränken:

\[
\int d^4 k \, \delta(k_0 - E(\mathbf{k})) = \int dk_0 \, d^3 k \, \delta(k_0 - E(\mathbf{k})).
\]

Die Integration über \(k_0\) mit der Delta-Funktion ergibt:

\[
\int dk_0 \, \delta(k_0 - E(\mathbf{k})) f(k_0) = f(E(\mathbf{k})).
\]

Somit reduziert sich das vierdimensionale Volumenelement auf:

\[
d^4 k = d^3 k \, \frac{1}{2 E(\mathbf{k})}.
\]

Die Transformation von \(dE\) zu \(dk_0\) erfordert die Berücksichtigung des Jacobians. Da \(k_0 = E(\mathbf{k})\), ergibt sich:

\[
dk_0 = \frac{\partial E}{\partial k_i} dk_i = \frac{k_i}{E} dk_i.
\]

Daher ist das dreidimensionale Integrationsmaß

\[
\frac{d^3 k}{2 E(\mathbf{k})}
\]

aus dem Lorentz-invarianten \(d^4 k\) abgeleitet und somit ebenfalls Lorentz-invariant. Die Normierung durch \((2 \pi)^3\) ist eine Konvention, die die Invarianz nicht beeinflusst.

%CHECK: Der Jacobian bei der Transformation von \(dE\) zu \(dk_0\) wurde korrekt behandelt. Die Rolle der Delta-Funktion bei der Integration über \(k_0\) wurde explizit erklärt.

\subsection*{(4.2) [v3]}
\paragraph{Aufgabentext.}
Use this to argue that $2 E(\mathbf{k}) \delta^{(3)}\left(\mathbf{k}-\mathbf{k}^{\prime}\right)$ is a Lorentz invariant distribution.
Hint: Start from the Lorentz-invariant expression $\frac{d^4 k}{(2 \pi)^3} \delta\left(k^2-m^2\right) \theta\left(k_0\right)$ and use $\delta\left(x^2-x_0^2\right)=\frac{1}{2\left|x_0\right|}\left(\delta\left(x+x_0\right)+\delta\left(x-x_0\right)\right)$.

\paragraph{Lösung.}
Um die Lorentz-Invarianz der Verteilung \(2 E(\mathbf{k}) \delta^{(3)}\left(\mathbf{k}-\mathbf{k}^{\prime}\right)\) zu zeigen, beginnen wir mit dem Lorentz-invarianten Ausdruck

\[
\frac{d^4 k}{(2 \pi)^3} \delta\left(k^2-m^2\right) \theta\left(k_0\right).
\]

Dieser Ausdruck ist Lorentz-invariant, da das Maß \(d^4 k\) und die Dirac-Delta-Funktion \(\delta\left(k^2-m^2\right)\) invariant sind, und die Heaviside-Funktion \(\theta\left(k_0\right)\) sicherstellt, dass nur positive Energien berücksichtigt werden.

Wir verwenden die Identität für die Delta-Funktion

\[
\delta\left(k^2-m^2\right) = \frac{1}{2|E(\mathbf{k})|} \left( \delta\left(k_0 - E(\mathbf{k})\right) + \delta\left(k_0 + E(\mathbf{k})\right) \right),
\]

wobei \(E(\mathbf{k}) = \sqrt{\mathbf{k}^2 + m^2}\) die positive Energie ist. Da \(\theta(k_0)\) nur den positiven Energieanteil berücksichtigt, bleibt nur der Term \(\delta\left(k_0 - E(\mathbf{k})\right)\) übrig. Somit reduziert sich der Ausdruck auf

\[
\frac{d^4 k}{(2 \pi)^3} \frac{1}{2 E(\mathbf{k})} \delta\left(k_0 - E(\mathbf{k})\right).
\]

Durch Integration über \(k_0\) erhalten wir

\[
\int dk_0 \frac{1}{2 E(\mathbf{k})} \delta(k_0 - E(\mathbf{k})) = \frac{1}{2 E(\mathbf{k})}.
\]

Dies führt zu

\[
\frac{d^3 \mathbf{k}}{(2 \pi)^3} \frac{1}{2 E(\mathbf{k})}.
\]

Um die Lorentz-Invarianz von \(2 E(\mathbf{k}) \delta^{(3)}\left(\mathbf{k}-\mathbf{k}^{\prime}\right)\) zu zeigen, multiplizieren wir mit \(2 E(\mathbf{k})\), was ergibt

\[
\frac{d^3 \mathbf{k}}{(2 \pi)^3} \delta^{(3)}\left(\mathbf{k}-\mathbf{k}^{\prime}\right).
\]

Nun zeigen wir die Invarianz von \(d^3 \mathbf{k}\) und \(\delta^{(3)}\left(\mathbf{k}-\mathbf{k}^{\prime}\right)\) unter Lorentz-Transformationen. Das Maß \(d^3 \mathbf{k}\) ist invariant unter räumlichen Drehungen, da es ein Skalarprodukt über den Raum ist. Für die Delta-Funktion gilt:

\[
\delta^{(3)}\left(\mathbf{k}-\mathbf{k}^{\prime}\right) = \prod_{i=1}^{3} \delta\left(k_i - k_i^{\prime}\right),
\]

welche unter Lorentz-Transformationen invariant bleibt, da die Transformationen linear sind und die Determinante der Jacobi-Matrix für Drehungen und Boosts gleich eins ist.

Daher ist der gesamte Ausdruck Lorentz-invariant, was zeigt, dass \(2 E(\mathbf{k}) \delta^{(3)}\left(\mathbf{k}-\mathbf{k}^{\prime}\right)\) eine Lorentz-invariante Verteilung ist.

%CHECK: Herleitung der Invarianz von \(d^3 \mathbf{k}\) und \(\delta^{(3)}\) unter Lorentz-Transformationen explizit gezeigt.

\end{document}\n